\section*{Capítulo \ref{ch:b2:population-genetics} --- Genética de poblaciones}

\begin{solution}{exo:b2:population-genetics:1}
$p(M) = (1280 + 320)/2000 = 0.8$, $q(N) = 0.2$. Esperados: $MM$,
$0.64\times 1000 = 640$; $MN$, $2\times 0.8\times 0.2\times 1000 =
320$; $NN$, 40 --- exactamente los números observados. La población está
en equilibrio en este locus.
\end{solution}

\begin{solution}{exo:b2:population-genetics:2}
$q = \sqrt{1/20000} = 0.0071$; portadores, $2pq \approx 0.014$, una
persona de cada 70. Los portadores superan a los afectados en una
proporción de $2p/q \approx
280$ a uno.
\end{solution}

\begin{solution}{exo:b2:population-genetics:3}
\omterm{def:b2:population-genetics:fitness}{Eficacia biológica}: el número esperado de descendientes de un \omterm{def:b2:meiosis-heredity:vocabulary}{genotipo}
en relación con el mejor. \omterm{def:b2:population-genetics:fitness}{Coeficiente de selección} $s$: el déficit de
eficacia, $w = 1 -
s$. Heredabilidad $h^{2}$: la fracción de la varianza fenotípica que es
genética aditiva, o, de forma equivalente, la pendiente de la regresión
de los descendientes sobre los progenitores. Fuerzas: la \omterm{def:b2:genome-diversification:mutation}{mutación}, la
selección, la deriva, la migración y el apareamiento no aleatorio (este
último cambia las frecuencias genotípicas, no las alélicas).
\end{solution}

\begin{solution}{exo:b2:population-genetics:4}
Direccional: el melanismo de la polilla del abedul, con las aves que cazan
a la vista sobre troncos ennegrecidos por el hollín como agente.
Estabilizadora: el peso al nacer en la especie humana, con la mortalidad
de los bebés muy pequeños y muy grandes como agente. Equilibradora: la
anemia falciforme, con el paludismo (contra $AA$) y la anemia (contra
$SS$) actuando juntos como agentes.
\end{solution}

\begin{solution}{exo:b2:population-genetics:5}
$q_t = q_0/(1 + tq_0)$: reducirla a la mitad lleva $1/q_0 = 50$
generaciones; llegar a $0.001$ lleva $1/0.001 - 1/0.02 = 950$
generaciones. Casi todas las copias del \omterm{def:b2:meiosis-heredity:vocabulary}{alelo} se encuentran en
\omterm{def:b2:meiosis-heredity:vocabulary}{heterocigotos}, invisibles para la selección; impedir que los afectados se
reproduzcan (cosa que la naturaleza ya hace, puesto que mueren) no cambia
nada medible en el tiempo de una vida humana.
\end{solution}

\begin{solution}{exo:b2:population-genetics:6}
$p = 0.3$: $\bar w = 0.09 + 0.42 + 0.8\times 0.49 = 0.902$; $p' =
(0.09 + 0.21)/0.902 = 0.333$; $\Delta p = +0.033$. $p = 0.9$: $\bar w
= 0.998$, $p' = 0.9018$, $\Delta p = +0.0018$. Con $p = 0.9$, solo
$q^{2} = 0.01$ de la población es $aa$: el \omterm{def:b2:meiosis-heredity:vocabulary}{alelo} contra el que actúa la
selección se esconde en los \omterm{def:b2:meiosis-heredity:vocabulary}{heterocigotos}, y $\Delta p = spq^{2}/\bar w$
lleva el factor $q^{2}$.
\end{solution}

\begin{solution}{exo:b2:population-genetics:7}
Con $s = 0.12$ contra $AA$ y $t = 0.8$ contra $SS$, $\hat q =
s/(s + t) = 0.12/0.92 = 0.13$. Nacimientos afectados: $\hat q^{2} = 0.017$,
alrededor de un niño de cada 60 --- el precio de la protección de la que
gozan los \omterm{def:b2:meiosis-heredity:vocabulary}{heterocigotos}.
\end{solution}

\begin{solution}{exo:b2:population-genetics:8}
Letal \omterm{def:b2:meiosis-heredity:vocabulary}{recesivo}: $\hat q = \sqrt{\mu} = \sqrt{2\times 10^{-5}} =
0.0045$; nacimientos afectados, $\hat q^{2} = \mu = 2\times 10^{-5}$, uno
de cada \num{50000}. Letal \omterm{def:b2:meiosis-heredity:vocabulary}{dominante}: cada copia se elimina en la
generación en que aparece, de modo que los nacimientos afectados son
simplemente las nuevas \omterm{def:b2:genome-diversification:mutation}{mutaciones}, $2\mu = 4\times 10^{-5}$ (dos gametos
por nacimiento), uno de cada \num{25000} --- el \omterm{def:b2:meiosis-heredity:vocabulary}{recesivo} esconde sus
\omterm{def:b2:meiosis-heredity:vocabulary}{alelos} en una proporción de doscientos a uno; el \omterm{def:b2:meiosis-heredity:vocabulary}{dominante} no esconde
nada.
\end{solution}

\begin{solution}{exo:b2:population-genetics:9}
$H_t = 0.5\,(1 - 1/100)^{t}$: al cabo de 10 generaciones, $0.45$; de 50,
$0.30$; de 200, $0.067$. Para $(1 - 1/2N)^{100} = 0.9$:
$100/(2N) \approx -\ln 0.9 = 0.105$, de modo que $N \approx 475$, unos 500
individuos reproductores.
\end{solution}

\begin{solution}{exo:b2:population-genetics:10}
Veinte copias génicas: $P(\text{ausente}) = 0.9^{20} = 0.12$. Una sola
copia entre 20 tiene una frecuencia de $0.05$, y un \omterm{def:b2:meiosis-heredity:vocabulary}{alelo} neutro se fija
con una probabilidad igual a su frecuencia: $0.05$. En un continente de,
por ejemplo, \num{100000} aves, una sola copia tiene una probabilidad de
$5\times 10^{-6}$ de fijarse. La isla es un lugar donde los \omterm{def:b2:meiosis-heredity:vocabulary}{alelos} raros
se pierden y donde los que sobreviven pueden imponerse: la deriva a la vez
empobrece y diferencia.
\end{solution}

\begin{solution}{exo:b2:population-genetics:11}
$R = h^{2}S = 0.3\times 1.4\times \qty{1000}{L} = \qty{420}{L}$ por
generación. Con toros del \qty{1}{\%} superior, el diferencial medio es de
$(1.4 + 2.7)/2 = 2.05$ desviaciones típicas y $R =
0.3\times 2050 = \qty{615}{L}$. La respuesta se frena porque la selección
agota la varianza aditiva --- los \omterm{def:b2:meiosis-heredity:vocabulary}{alelos} favorables llegan a la fijación y
$h^{2}$ disminuye --- y porque los cambios correlacionados de otros rasgos
(fertilidad, salud) imponen costes que la selección sobre la producción
sola no ve.
\end{solution}

\begin{solution}{exo:b2:population-genetics:12}
La selección ve \omterm{def:b2:meiosis-heredity:vocabulary}{fenotipos}. Un \omterm{def:b2:meiosis-heredity:vocabulary}{alelo} \omterm{def:b2:meiosis-heredity:vocabulary}{recesivo} en un \omterm{def:b2:meiosis-heredity:vocabulary}{heterocigoto} no produce
\omterm{def:b2:meiosis-heredity:vocabulary}{fenotipo}, de modo que la selección no puede eliminarlo, y su frecuencia
baja como $1/t$, cada vez más despacio. Una \omterm{def:b2:genome-diversification:mutation}{mutación} neutra no produce
ninguna diferencia de eficacia; su destino depende solo de la deriva, y se
fija con probabilidad $1/2N$. Un carácter de heredabilidad nula puede
estar sometido a una selección intensa --- los supervivientes pueden
diferir mucho de la población ---, y sin embargo la generación siguiente
no cambia, porque las diferencias no se heredaban. La selección solo puede
actuar sobre diferencias heredables que suponen una diferencia de
eficacia; todo lo demás evoluciona por azar o no evoluciona.
\end{solution}

\begin{solution}{pb:b2:population-genetics:1}
\textbf{1.} Fracción melánica: $1 - q^{2} = 1 - 0.999^{2} = 0.002$,
una polilla de cada 500; la fracción de \omterm{def:b2:meiosis-heredity:vocabulary}{alelos} $M$ en \omterm{def:b2:meiosis-heredity:vocabulary}{heterocigotos} es
$2pq/(2pq + 2p^{2}) = q = 0.999$: prácticamente todos.
\textbf{2.} Con $\bar w = 1 - sq^{2}$, las copias del \omterm{def:b2:meiosis-heredity:vocabulary}{alelo} claro tras la
selección son $pq\cdot 1 + q^{2}(1 - s)$ sobre $\bar w$, de modo que $q'
= (pq + q^{2} - sq^{2})/(1 - sq^{2}) = q(1 - sq)/(1 - sq^{2})$.
\textbf{3.} $q_1 = 0.999\times(1 - 0.2997)/(1 - 0.2994) = 0.99857$:
un cambio de $-0.0014$.
\textbf{4.} $q = \sqrt{0.1} = 0.32$.
\textbf{5.} La recurrencia, iterada, lleva 28 generaciones para $s =
0.3$; la aproximación da el mismo orden (el cambio es de
$-0.0003$ al principio y después se acelera a medida que crece $p$). Las
50 generaciones observadas corresponden a un $s$ menor, de alrededor de
$0.2$ (la recurrencia da 45 generaciones para $s = 0.2$): una selección
del \qtyrange{20}{30}{\%} por generación, enorme para lo que es habitual en
la mayoría de los \omterm{def:b2:meiosis-heredity:vocabulary}{alelos}.
\textbf{6.} Con $w_{MM} = w_{Mm} = 1 - s$, $w_{mm} = 1$: $\bar w = 1
- s(1 - q^{2})$ y $p' = p(1 - s)/\bar w$, de modo que $\Delta p = p[(1 - s) -
\bar w]/\bar w = -spq^{2}/\bar w$. La forma clara solo está favorecida
como \omterm{def:b2:meiosis-heredity:vocabulary}{homocigoto} $mm$, una fracción $q^{2} = 0.1$ al principio: el regreso
empieza despacio. Cuando $q \to 1$, $\Delta p \approx -sp$ y $M$
disminuye de forma exponencial, en un factor $1 - s$ por generación: un
\omterm{def:b2:meiosis-heredity:vocabulary}{alelo} \omterm{def:b2:meiosis-heredity:vocabulary}{dominante} no tiene dónde esconderse. El ascenso fue la imagen
especular:
$\Delta p \approx +sp$ mientras $M$ era raro (comienzo exponencial), y
después
$\Delta q \approx -sq^{2}p \to$ un avance lentísimo a medida que el \omterm{def:b2:meiosis-heredity:vocabulary}{alelo}
claro se escondía en los \omterm{def:b2:meiosis-heredity:vocabulary}{heterocigotos}. La recurrencia da 27 generaciones
para el regreso de $p = 0.68$ a $0.001$, 16 de ellas para bajar a una
\omterm{def:b2:population-genetics:frequencies}{frecuencia alélica} del \qty{5}{\%}.
\textbf{7.} $24$ copias génicas: $0.95^{24} = 0.29$.
\textbf{8.} $H_1 = 0.4\,(1 - 1/24) = 0.383$.
\textbf{9.} $H_{100} = 0.383\,(1 - 1/80)^{100} = 0.383\times 0.284 =
0.109$: el \qty{27}{\%} de los $0.4$ del continente.
\textbf{10.} $P(\text{fijación}) = 1/(2N) = 1/80 = 0.0125$; tiempo
esperado hasta la fijación, $4N = 160$ generaciones.
\textbf{11.} Sin endogamia, $q^{2} = 0.0025$; con $F = 0.3$,
$q^{2} + Fpq = 0.0025 + 0.3\times 0.95\times 0.05 = 0.017$: siete
veces más aves afectadas.
\textbf{12.} Un migrante por generación ($m = 1/40$) basta para impedir
que la isla diverja en los loci neutros y para reponer los \omterm{def:b2:meiosis-heredity:vocabulary}{alelos} que la
deriva ha perdido; también importa los \omterm{def:b2:meiosis-heredity:vocabulary}{alelos} perjudiciales del continente
a su frecuencia continental, que, con la endogamia de la isla, quedan
expuestos en \omterm{def:b2:meiosis-heredity:vocabulary}{homocigotos} --- la carga la fija el equilibrio entre la
inmigración y la exposición.
\textbf{13.} $\hat q = \sqrt{10^{-5}} = 0.0032$; nacimientos afectados, $\hat
q^{2} = 10^{-5}$.
\textbf{14.} $\hat q = \sqrt{10^{-5}/0.1} = 0.01$; nacimientos afectados,
$10^{-4}$: una selección diez veces más débil permite que el \omterm{def:b2:meiosis-heredity:vocabulary}{alelo} sea tres
veces más común y la enfermedad diez veces más.
\textbf{15.} $\hat p = \mu/(hs) = 10^{-5}/0.05 = 2\times 10^{-4}$;
portadores, $2\hat p = 4\times 10^{-4}$.
\textbf{16.} De $\hat q = 0.0032$ a $0.01$: el \omterm{def:b2:meiosis-heredity:vocabulary}{alelo} sube a un ritmo
fijado por la \omterm{def:b2:genome-diversification:mutation}{mutación}, $\mu = 10^{-5}$ por generación, de modo que la
aproximación lleva del orden de $\Delta q/\mu \approx 700$ generaciones ---
\num{20000} años: cualquier decisión tomada hoy la notará una población
que no se parecerá en nada a la nuestra.
\textbf{17.} Cada locus aporta $2\hat q = 0.0063$ copias letales por
persona; sobre \num{20000} loci, unas $130$ --- pero eso supone que todos
los genes pueden mutar a un letal \omterm{def:b2:meiosis-heredity:vocabulary}{recesivo}; con la estimación habitual de
unos pocos miles de genes esenciales, una persona porta en estado
\omterm{def:b2:meiosis-heredity:vocabulary}{heterocigoto} desde varios hasta unas pocas decenas de equivalentes letales.
\textbf{18.} Cada uno está en un locus distinto y cada uno es raro, de
modo que la probabilidad de que los dos miembros de una pareja tomada al
azar porten el mismo es pequeña; los primos comparten una octava parte de
sus genes por ascendencia, de modo que, para cada uno de los varios letales
de un primo, el riesgo de que el hijo sea \omterm{def:b2:meiosis-heredity:vocabulary}{homocigoto} es de
$1/16$ por letal compartido --- el exceso de mortalidad infantil de los
matrimonios entre primos medido es de unos pocos puntos porcentuales.
\textbf{19.} $S = 1.16\times \qty{800}{L} = \qty{930}{L}$; $R = 0.4\times
930 = \qty{370}{L}$.
\textbf{20.} Cinco generaciones: unos \qty{1860}{L}; menos en la práctica,
porque la varianza aditiva se agota, porque $h^{2}$ se estimó en el rebaño
original y porque el ambiente (alimentación, instalaciones) que produjo
las mejores vacas puede no transmitirse.
\textbf{21.} Diferencial medio de $(1.16 + 2.4)/2 = 1.78$ desviaciones
típicas, $S = \qty{1424}{L}$, $R = 0.4\times 1424 = \qty{570}{L}$.
\textbf{22.} Los \omterm{def:b2:asexual-reproduction:clone}{clones} no tienen varianza genética, de modo que $h^{2} = 0$ y $R =
0$ sea cual sea $S$: la heredabilidad es una propiedad de una población
--- la parte de su varianza que es genética ---, no una propiedad del
rasgo.
\textbf{23.} $R = 0.7\times 0.6 = \qty{0.42}{mm}$ más profundo.
\textbf{24.} La misma respuesta de $\qty{370}{L}$ exige $S = 370/0.1 =
\qty{3700}{L}$, es decir, $4.6$ desviaciones típicas --- conservar menos de
una vaca de cada diez mil: impracticable, y por eso los rasgos de baja
heredabilidad se mejoran despacio o por otros medios (pruebas de
descendencia, familias más numerosas).
\textbf{25.} Polilla: $s \approx 0.2$ a $0.3$, 28 generaciones con $0.3$
y 45 con $0.2$, frente a las 50 observadas; isla: $H_{100} \approx 0.11$,
el \qty{27}{\%} de la del continente; rebaño: \qty{370}{L} por generación
seleccionando solo las vacas, \qty{570}{L} con toros seleccionados.
\end{solution}
